\chapter*{Introduction g\'en\'erale}
\addcontentsline{toc}{chapter}{Introduction g\'en\'erale}

\guided{Cette introduction doit pr\'esenter clairement : le contexte g\'en\'eral, la probl\'ematique, les objectifs, le p\'erim\`etre, l'organisme d'accueil si pertinent, et l'organisation du rapport. Longueur recommand\'ee : 1 \`a 2 pages. Pas de figures ni de tableaux dans cette partie.}

Dans un contexte de transformation num\'erique et industrielle rapide, les projets de fin d'\'etudes constituent une situation d'apprentissage sommative o\`u l'\'el\`eve-ing\'enieur doit d\'emontrer sa capacit\'e \`a analyser un probl\`eme complexe, \`a proposer une d\'emarche de r\'esolution rigoureuse et \`a valider une solution de mani\`ere argument\'ee.

Ce travail s'inscrit dans le cadre d'un projet r\'ealis\'e au sein de \CompanyName. Le sujet porte sur \todo{pr\'eciser le sujet en une phrase claire}. La probl\'ematique principale peut \^etre formul\'ee comme suit : \todo{formuler la question centrale du projet}. Les objectifs retenus visent \`a \todo{lister les objectifs sp\'ecifiques en coh\'erence avec le sujet valid\'e}.

Le p\'erim\`etre du projet couvre \todo{indiquer ce qui est inclus et ce qui est hors p\'erim\`etre}. Cette pr\'ecision est indispensable pour \^eviter une pr\'esentation descriptive sans logique de cadrage. \disciplinehint

Le pr\'esent rapport est organis\'e comme suit. Le premier chapitre expose l'\'etat de l'art et le positionnement du projet. Le deuxi\`eme chapitre pr\'esente l'analyse des besoins, les sp\'ecifications et la m\'ethodologie. Le troisi\`eme chapitre d\'etaille la conception et l'architecture. Le quatri\`eme chapitre traite la r\'ealisation ou l'exp\'erimentation. Le cinqui\`eme chapitre expose les tests, la validation et la discussion critique des r\'esultats. Enfin, la conclusion g\'en\'erale synth\'etise les apports du travail et propose des perspectives.
